\section*{Chapter \ref{ch:b2:specification} --- Specification}

\begin{solution}{exo:b2:specification:1}
% figdata: attenuation
At $\sigma_e^2 = 1$ the estimate is 0.99, against a limit of 1; the \omterm{def:b2:specification:reliability}{reliability ratio} is $1/(1 + 1) = 0.5$.
\end{solution}

\begin{solution}{exo:b2:specification:2}
% code: listings/engel.py
A 1 per cent rise in income goes with 0.856 per cent more food spending. A 10 per cent rise raises food spending by about
8.6 per cent (exactly $1.1^{0.856} - 1 = 8.5$ per cent), less than income, so the budget share falls.
\end{solution}

\begin{solution}{exo:b2:specification:3}
% code: listings/dummies.py
With a constant, a fifth dummy would make the dummies sum to the constant: the \omterm{def:b2:specification:dummy-trap}{dummy-variable trap}. Size 1 is the omitted
category. A household of three: $e^{0.440} - 1 = 55$ per cent more.
\end{solution}

\begin{solution}{exo:b2:specification:4}
If none of the twenty mattered, at least one would be significant at 5 per cent with probability about 0.64. One
significant result in twenty is what chance alone predicts; the reader should ask for all twenty results, or a threshold
of $0.05/20$.
\end{solution}

\begin{solution}{exo:b2:specification:5}
% code: exercises/solutions.py percent
About 15 per cent; exactly $e^{0.15} - 1 = 16.2$ per cent.
\end{solution}

\begin{solution}{exo:b2:specification:6}
$\lambda\beta = 2/(1 + \sigma_e^2)$: 1.6 needs $\sigma_e^2 = 0.25$, 1 needs $\sigma_e^2 = 1$; the listing prints 1.605 and
1.006.
\end{solution}

\begin{solution}{exo:b2:specification:7}
% code: exercises/solutions.py engel_quadratic
The coefficients are 0.872 on centred log income and $-0.077$ on its square ($p = 0.014$): the elasticity falls with
income, $0.872 - 2 \times 0.077 \times (\ln\mathit{inc} - \overline{\ln\mathit{inc}})$. It is 0.96 at the 10th percentile of
income and 0.79 at the 90th.
\end{solution}

\begin{solution}{exo:b2:specification:8}
% code: exercises/solutions.py proxy_coef; tests/test_solutions.py proxy_plim
A better proxy removes more of the bias: 0.083, against 0.091 with the noisier score and 0.080 with ability itself. The
population value is $0.08 + 0.10\,\delta$, with $\delta = (1.5 \times 1.09 - 1.5)/(6.25 \times 1.09 - 1.5^2) = 0.030$.
\end{solution}

\begin{solution}{exo:b2:specification:9}
The programme may change the worker's occupation; comparing trained and untrained workers in the same occupation removes
that channel and compares groups that differ in unobserved ways. Good controls are fixed before the programme: age,
schooling, earlier earnings and occupation before the programme.
\end{solution}

\begin{solution}{exo:b2:specification:10}
By the theorem the slope on $x_1$ is the slope of $y$ on $\tilde x_1$, the residual of $x_1$ on $x_2$. Since $e$ is
uncorrelated with $x_2$, the regression of $x_1$ on $x_2$ has the same limit as that of $x_1^*$, so $\tilde x_1 = \tilde
x_1^* + e$ in the limit, and $y$ moves with
$\tilde x_1^*$ at rate $\beta_1$; so the slope converges to $\beta_1\operatorname{Var}(\tilde x_1^*)/(\operatorname{Var}(\tilde x_1^*) + \operatorname{Var}(e)) = \lambda_1\beta_1$.
Since $\operatorname{Var}(\tilde x_1^*) \leq \operatorname{Var}(x_1^*)$ and the noise variance is unchanged, $\lambda_1 \leq
\lambda$.
\end{solution}

\begin{solution}{exo:b2:specification:11}
If $y = y^* + v$ with $v$ independent of $x$, then $y = \beta_0 + \beta x + (u + v)$ and $x$ is still uncorrelated with the
error: the slope is consistent. The error variance rises from $\sigma_u^2$ to $\sigma_u^2 + \sigma_v^2$, and the variance
of the slope rises in proportion.
\end{solution}

\begin{solution}{exo:b2:specification:12}
% code: exercises/solutions.py any_rejection
The probability of a union is at most the sum of the probabilities: $k\alpha$. Under independence the probability that
none rejects is $(1 - \alpha)^k$. With $\alpha = 0.05$: bound 1 (useless), exact 0.642. With $\alpha = 0.0025$: bound
0.05, exact 0.0488.
\end{solution}

\begin{solution}{pb:b2:specification:1}
% code: exercises/solutions.py brenland_food, corrected, bonferroni_z, any_rejection; listings/dummies.py
\textbf{Part I.}
1. 0.288; $R^2 = 0.181$.
2. A necessity: the elasticity is below one.
3. By about 2.9 per cent.
4. It falls, since food spending rises less than in proportion.
5. The slope is then an elasticity, free of units, and the relation is closer to linear.
\textbf{Part II.}
6. 0.196; $R^2 = 0.924$.
7. Larger households have both more income and more mouths to feed; without size, the income coefficient picked up part
of the effect of size.
8. 0.246, 0.440, 0.597 and 0.738 for sizes 2 to 5; a household of five spends $e^{0.738} - 1 = 109$ per cent more than a
single person with the same income.
9. Size 1 is the omitted category; a dummy for it with the constant would be the \omterm{def:b2:specification:dummy-trap}{dummy-variable trap}.
10. $+0.001$ (s.e.\ 0.004): owners and tenants of the same size and income spend the same on food.
11. Tenants 0.187, owners 0.211.
12. Yes: $t = 2.64$ (from the unrounded coefficient and standard error), above 1.96.
\textbf{Part III.}
13. Toward zero.
14. $0.196/0.8 = 0.245$.
15. By the \omterm{def:b2:multiple-regression:fwl}{Frisch--Waugh--Lovell theorem} the elasticity comes from the part of income not predicted by size, and it is the
reliability of that part which sets the attenuation.
16. No: without size, the reliability of the whole of reported income would apply, which is higher, and the elasticity
itself would be a different one.
17. No: error in the outcome raises the standard errors but leaves the slope unbiased.
18. Then the error is negatively correlated with true income and the classical correction is wrong; the bias depends on
how reporting varies with income.
\textbf{Part IV.}
19. $1 - 0.95^{20} = 0.642$.
20. $0.05/20 = 0.0025$; $|t| > 3.02$.
21. No: its $t$ is 2.64.
22. The regressors, the form and the sample of the regression it will report, and how many others it will try.
23. It was chosen for a reason given in advance (larger households eat more), not found by the search.
24. No: choosing after seeing the data is a search too, and the number of tries is then not twenty but unknown.
25. Brenland's food--income elasticity, holding household size fixed, is 0.196 as measured and $0.196/0.8 = 0.245$
corrected for the reliability of reported income; with twenty tries, a result should count only if its $p$-value is below
0.0025.
\end{solution}
